\begin{textsample}{POS Dim 4 – 2015\_c5 – Score 1.00 – 2015\_c5/enade\_2015\_jorn\_aec7ff8cf3.txt}  \label{ex:f4_pos_165}
CAPÍTULO 5 DISTRIBUIÇÃO DOS CONCEITOS 5.1 PANORAMA NACIONAL DA DISTRIBUIÇÃO DOS CONCEITOS O Conceito Enade é um indicador de qualidade que avalia o desempenho dos estudantes a partir dos resultados do Enade.
Os valores possíveis vão de 1 ( pior situação ) a 5 ( melhor situação ).
As Unidades de \textbf{Observação} com apenas um ou sem nenhum concluinte participante não obtêm o Conceito Enade, ficando Sem Conceito ( SC ).
A Tabela 5.1 apresenta a quantidade e distribuição de cursos de Jornalismo participantes do Enade/2015, por faixa de conceito e Grande Região.
Mais uma vez, ressaltando que a diferença entre os cursos tabulados neste capítulo e no capítulo 2 corresponde aos cursos sem conceito, em princípio, aqueles sem alunos concluintes que participassem da prova ou que tivessem tão somente um aluno concluinte18.
Observando -se os dados da Tabela 5.1, nota -se que, dos 275 cursos participantes, 116 ( 42,2 \% ) classificaram -se com conceito 3, o valor modal.
Este foi também o conceito modal em todas as demais regiões : Norte ( 40,9 \% ), Nordeste ( 46,9 \% ), Sudeste ( 39,7 \% ), Sul ( 40,0 \% ) e Centro-Oeste ( 52,2 \% ).
O conceito 4 foi o segundo mais frequente em nível nacional ( 23,3 \%, correspondendo a 64 cursos ) e o conceito 2, o terceiro ( 21,5 \%, correspondendo a 59 cursos ).
Houve, ainda, 18 cursos ( 6,5 \% ) que receberam conceito 5 e sete cursos ( 2,5 \% ) que receberam conceito 1.
Dos 275 cursos de Jornalismo, 11 cursos ( 4,0 \% ) ficaram sem conceito ( SC ). 18 Estes cursos com somente um concluinte não foram considerados no capítulo 2.
A região Norte participou com 22 cursos ou 8,0 \% do total nacional.
Nesta região, o conceito 3 foi atribuído a nove cursos dos 22 participantes, o que equivale a 40,9 \% do total regional, além de ser o conceito modal na região.
Apenas um curso ficou sem conceito ( SC ) ou com conceito 1.
Oito cursos foram avaliados com conceito 2 e três cursos obtiveram conceito 4, correspondendo, respectivamente, a 36,4 \% e 13,6 \% da região.
Nenhum curso foi avaliado com conceito 5 nessa região.
A região Nordeste participou com 49 cursos ou 17,8 \% do total nacional.
Nesta região, 23 cursos, 46,9 \% em termos regionais, obtiveram conceito 3, o conceito modal para a região.
Os conceitos 1 e 2 foram atribuídos, respectivamente, a um e 12 cursos ( 2,0 \% e 24,5 \% ).
Dez cursos obtiveram o conceito 4 ( 20,4 \% ) e dois cursos receberam conceito 5 ( 4,1 \% ).
Apenas um curso ficou sem conceito ( SC ) nessa região.
Dos 126 cursos participantes da região Sudeste ( 45,8 \% do total nacional ), o conceito 3 também foi o conceito modal, e foi obtido por 50 cursos ( 39,7 \% ).
O conceito 1 foi atribuído a quatro cursos ( 3,2 \% ), o conceito 2 a 29 cursos ( 23,0 \% ), o conceito 4 a 28 cursos ( 22,2 \% ) e o 5 a nove cursos, correspondendo a 7,1 \%.
Seis cursos ( 4,8 \% ) ficaram sem conceito ( SC ).
Os 55 cursos da região Sul ficaram distribuídos em todas as faixas de conceitos, à exceção do conceito 1 que não foi atribuído a nenhum curso.
Há predominância do conceito 3 com 40,0 \% ( conceito modal ), ou seja, 22 dos 55 cursos participantes na região Sul.
O conceito 2 foi atribuído a seis cursos ( 10,9 \% ) e o conceito 4, a 19 cursos ( 34,5 \% ).
O conceito 5 foi recebido por seis cursos ( 10,9 \% ).
Apenas dois cursos na região Sul ficaram sem conceito ( SC ).
Já dos 23 cursos participantes na região Centro-Oeste, 12 receberam conceito 3, o conceito modal ( 52,2 \% ).
Nesta região, apenas um curso ficou sem conceito ( SC ) e os conceitos 1 e 5, também foram atribuídos a apenas um curso cada ( 4,3 \% cada ).
Já o conceito 2 foi atribuído a quatro cursos, correspondendo a 17,4 \% dos cursos nessa região.
O Gráfico 5.1 apresenta a distribuição do Conceito Enade segundo Grande Região : poligonais mais à esquerda representam Grandes Regiões com pior distribuição, e poligonais mais à direita, Grandes Regiões com melhores conceitos.
A região Norte ( linha azul ) apresenta os cursos em IES com conceitos mais baixos, ao passo que a região Sul ( linha roxa ) apresenta a distribuição com valores maiores.
Cursos em IES das regiões Sudeste ( linha vermelha ), Nordeste ( linha verde ) e Centro-Oeste ( linha preta ) apresentam uma situação intermediária. 5.2 CONCEITOS POR CATEGORIA ADMINISTRATIVA E POR GRANDE REGIÃO A Tabela 5.2 apresenta a distribuição dos cursos participantes do Enade/2015 de Jornalismo, por Categoria Administrativa, de acordo com os conceitos por eles alcançados, segundo as Grandes Regiões brasileiras.
Dos 275 cursos participantes, 63 ( 22,9 \% ) eram ministrados em Instituições Públicas e 212 ( 77,1 \% ), em Instituições Privadas.
De acordo com as informações da Tabela 5.2, em termos nacionais, entre Instituições Públicas e Privadas, nota -se que das 18 IES que receberam o conceito 5, 11 eram Públicas e sete Privadas.
Dos 63 cursos participantes de IES Públicas, o conceito 3 foi o valor modal, atribuído a 20 cursos ( 31,7 \% ).
Entre os demais cursos de Instituições Públicas participantes, dois cursos ( 3,2 \% ) foram avaliados com conceito 1, dez cursos ( 15,9 \% da categoria ) obtiveram conceito 2, 19 cursos ( 30,2 \% ) receberam conceito 4 e, como já comentado, 11 cursos ( 17,5 \% ) foram alocados ao conceito 5.
Nesta categoria, apenas um curso ficou sem conceito ( SC ).
Na rede privada, o conceito modal também foi 3, atribuído a 96 cursos dos 212 da categoria.
Entre os demais cursos participantes, cinco receberam conceito 1, e 49, conceito 2.
O conceito 4 foi atribuído a 45 cursos e o conceito 5, a sete cursos.
Nesta Categoria Administrativa, dez cursos ficaram sem conceito ( SC ).
Na análise por região, observa -se que, na região Norte, as Instituições Privadas participaram com 13 cursos ( 59,1 \% do total regional ), dos quais apenas um curso ficou sem conceito ( SC ) ou com conceitos 1 ou 4.
O conceito modal para as Instituições Privadas na região foi 3, atribuído a seis cursos, correspondendo a 46,2 \%, e o conceito As Instituições Privadas participaram com 13 cursos ( 59,1 \% do total regional ), dos quais apenas um curso ficou sem conceito ( SC ) ou com conceitos 1 ou 4.
O conceito modal para as Instituições Privadas na região foi 3, atribuído a seis cursos, correspondendo a 46,2 \%, e o conceito 2 foi atribuído a quatro cursos ( 30,8 \% ).
Nenhum curso foi avaliado com conceito 5.
As Instituições Públicas participaram com nove cursos na região Norte ( 40,9 \% do total regional ), dos quais quatro obtiveram conceito 2, o valor modal, seguido por três cursos ( 33,3 \% ) que receberam conceito 3 e dois cursos ( 22,2 \% ) que ficaram com conceito 4.
Nenhum curso ficou sem conceito ( SC ) ou recebeu os conceitos 1 ou 5 nesta combinação de categoria e região.
Na região Nordeste, a rede privada concentrou 30 dos 49 cursos participantes, o equivalente a 61,2 \% do total da região.
Dentre os cursos oferecidos por IES Privadas no Nordeste, um curso ficou alocado no conceito 1, dez cursos obtiveram conceito 2, 15 cursos ficaram com conceito 3, conceito modal, e quatro cursos foram avaliados com o conceito 4.
Nenhum curso ficou sem conceito ( SC ) ou com conceito 5.
As Instituições Públicas dessa região participaram com 19 cursos ( 38,8 \% ), dos quais oito ( 42,1 \% da categoria na região ) obtiveram conceito 3, o valor modal, seguido por seis cursos que receberam conceito 4 ( 31,6 \% ).
Dois cursos foram alocados no conceito 2 ( 10,5 \% ) e outros dois ( 10,5 \% ) no conceito 5.
Apenas um curso ficou sem conceito ( SC ) e nenhum obteve conceito 1 nesta combinação de categoria e região.
Na região Sudeste, a proporção de cursos da rede privada, 85,7 \%, foi mais elevada do que nas demais regiões brasileiras, correspondendo a 108 dos 126 cursos participantes.
Nesta categoria, na região Sudeste, o conceito modal foi 3 ( 47 cursos ).
Nesta combinação de Categoria Administrativa e Grande Região, seis cursos não receberam conceito ( SC ).
Os demais foram avaliados com conceito 1 ( dois cursos ), conceito 2 ( 27 cursos ), conceito 4 ( 21 cursos ) e conceito 5 ( cinco cursos ).
Entre os 18 cursos oferecidos por Instituições Públicas na região Sudeste, o conceito modal foi 4 ( sete cursos ), e os demais cursos receberam os conceitos 1 ( dois cursos ), 2 ( dois cursos ), 3 ( três cursos ) e 5 ( quatro cursos ).
Na região Sudeste, nesta categoria, nenhum curso ficou sem conceito ( SC ).
As Instituições Privadas concentraram 44 dos 55 cursos participantes da região Sul, 80,0 \% do total regional.
Desses, 18 ficaram com conceito 3, o conceito modal.
Nesta combinação de Categoria Administrativa e Grande Região, 16 cursos ficaram com conceito 4, seis cursos obtiveram conceito 2, enquanto dois cursos receberam conceito 5.
Nenhum curso foi associado ao conceito 1, enquanto dois cursos ficaram sem conceito ( SC ).
As Instituições Públicas na região Sul participaram com 11 cursos ( 20 \% ), que apresentaram como conceitos modais os conceitos 3 e 5, atribuídos, cada qual, a 4 cursos.
Três cursos foram avaliados com conceito 4 e nenhum recebeu os conceitos 1, 2 ou ficou sem conceito ( SC ).
Na região Centro-Oeste, 17 dos 23 cursos participantes eram de Instituições Privadas ( 73,9 \% em termos regionais ).
Destes, dez concentraram -se no conceito 3, conceito modal.
Dentre os demais, um recebeu conceito 1, dois obtiveram conceito 2 e três cursos foram avaliados com conceito 4.
Apenas um curso ficou sem conceito ( SC ) e nenhum recebeu conceito 5 nesta categoria.
A distribuição dos seis cursos oferecidos por Instituições Públicas apresentou uma moda distribuída em dois conceitos contíguos ( conceitos 2 e 3 ) e alocando 2 cursos em cada um deles.
Os conceitos 4 e 5 foram conferidos a um curso cada.
Nesta região, nenhum curso de IES pública ficou sem conceito ( SC ) ou recebeu conceito 1.
O Gráfico 5.2 apresenta a distribuição do Conceito Enade segundo a Categoria Administrativa da IES.
Os cursos em IES Públicas ( linha azul ) apresentam uma poligonal mais à direita do que os em IES Privadas ( linha verde ) e, consequentemente, uma distribuição de Conceitos Enade com valores maiores. 5.3 CONCEITOS POR ORGANIZAÇÃO ACADÊMICA E POR GRANDE REGIÃO Na Tabela 5.3, encontra -se a distribuição dos conceitos atribuídos aos cursos participantes do Enade/2015 na Área de Jornalismo, por Organização Acadêmica, segundo as Grandes Regiões brasileiras.
Dos 275 cursos de Jornalismo participantes, 146 eram oferecidos em Universidades ; 55 cursos em Centros Universitários ; e 74 cursos em Faculdades.
Esta distribuição corresponde, respectivamente, a 53,1 \%, 20,0 \% e 26,9 \% dos cursos.
De acordo com os dados apresentados, de todos os 18 cursos avaliados com conceito 5, 13 eram vinculados a Universidades.
Cinco cursos oferecidos em Universidade ficaram sem conceito ( SC ).
Esse tipo de Organização Acadêmica teve o conceito 3 como modal, atribuído a 65 cursos ( 44,5 \% ).
Os demais cursos de Universidades avaliados receberam os conceitos 1 ( dois cursos ), 2 ( 21 cursos ), 4 ( 40 cursos ) e conceito 5 ( 13 cursos, como já mencionado ).
Entre os cursos em Centros Universitários, o conceito modal foi 3, com 26 cursos ( 47,3 \% ).
Neste tipo de Organização Acadêmica, dois cursos ficaram sem conceito ( SC ), e os outros cursos receberam os conceitos 1 ( dois cursos ), 2 ( 16 cursos ), e 4 ( nove cursos ).
Nenhum curso recebeu conceito 5.
Nas Faculdades, quatro dos 74 cursos ficaram sem conceito ( SC ), e 25 ( 33,8 \% ) ficaram alocados no conceito 3, conceito modal.
Os demais cursos neste tipo de Organização Acadêmica receberam os conceitos 1 ( três cursos ), 2 ( 22 cursos ), 4 ( 15 cursos ) e 5 ( cinco cursos ).
Considerando -se separadamente as regiões brasileiras, verifica -se que, na região Norte, as Universidades concentraram dez dos 22 cursos participantes.
Quatro cursos ficaram alocados no conceito 3, conceito modal, seguidos de três cursos que receberam conceito 2, dois cursos que obtiveram conceito 4 e um curso que obteve conceito 1.
Nenhum curso ficou sem conceito ( SC ) ou recebeu conceito 5 nesta categoria e região.
Os Centros Universitários da região Norte foram representados por quatro cursos, sendo que dois cursos obtiveram conceito 2, conceito modal, e um curso recebeu o conceito 3.
Nenhum dos cursos oferecidos por Centros Universitários na região Norte ficou com o conceito 1, 4 ou 5 e apenas um curso ficou sem conceito ( SC ).
As Faculdades participaram com oito cursos na região Norte, sendo que esses receberam conceitos 2 ( três cursos ), 3 ( quatro cursos, conceito modal ) e 4 ( um curso ).
Nenhum curso na região Norte ficou sem conceito ou foi avaliado com conceitos 1 ou 5.
Na região Nordeste, as Universidades participaram com 24 dos 49 cursos na Área de Jornalismo.
O conceito modal foi 3, atribuído a 12 cursos, enquanto o conceito 4 foi atribuído a seis cursos, o 5 a dois cursos, e o conceito 2, a três cursos.
Apenas um dos cursos oferecidos por Universidades no Nordeste ficou sem conceito ( SC ) enquanto nenhum obteve conceito 1.
Os Centros Universitários contaram com sete cursos participantes na região Nordeste, sendo que um curso obteve conceito 2, quatro cursos tiraram conceito 3, conceito modal, e dois cursos foram avaliados com o conceito 4.
Nenhum curso ficou sem conceito ( SC ) ou obteve conceito 1 ou 5.
As Faculdades foram representadas por 18 cursos na região Nordeste, dos quais um tirou conceito 1, oito obtiveram conceito 2, conceito modal, sete ficaram com conceito 3 e dois foram avaliados com o conceito 4.
Nenhum curso ficou sem conceito ( SC ) ou obteve conceito 5.
Na região Sudeste, as Universidades concentraram 67 dos 126 cursos da região.
Entre os cursos em Universidades nessa região, o conceito modal foi 3, com 28 cursos, e três cursos ficaram sem conceito ( SC ).
Os demais cursos receberam os conceitos 1 ( um curso ), 2 ( 11 cursos ), 4 ( 19 cursos ) e 5 ( cinco cursos ).
Os Centros Universitários participaram com 29 cursos na região Sudeste, dos quais 14 obtiveram o conceito modal 3, enquanto apenas um curso ficou sem conceito ( SC ).
Nenhum curso recebeu conceito 5 nessa categoria e região.
Os demais receberam os conceitos 1 ( um curso ), 2 ( dez cursos ) e 4 ( três cursos ).
As Faculdades foram representadas por 30 cursos na região Sudeste, dos quais apenas dois cursos ficaram sem conceito ( SC ) e dois cursos tiraram o conceito 1.
A distribuição dos cursos oferecidos por Faculdades na região Sudeste apresentou uma moda distribuída em dois conceitos contíguos, alocando 8 cursos em cada um dos conceitos 2 e 3.
Os demais cursos receberam os conceitos 4 ( seis cursos ) e 5 ( quatro cursos ).
Dos 55 cursos da região Sul, 33 eram de Universidades, para os quais o conceito modal foi 3, atribuído a 14 cursos.
Nesse tipo de organização, apenas um curso ficou sem conceito ( SC ), e nenhum recebeu o conceito 1.
Os demais receberam os conceitos 2 ( dois cursos ), 3, conceito modal, ( 14 cursos ), 4 ( 11 cursos ) e 5 ( cinco cursos ).
Dos 11 cursos participantes de Centros Universitários da região Sul, três receberam conceito 2, seis obtiveram conceito 3, conceito modal, e dois cursos foram avaliados com o conceito 4.
Nenhum curso ficou sem conceito ( SC ) ou recebeu conceito 1 ou 5.
Foram 11 os cursos vinculados a Faculdades na região Sul, sendo que seis desses receberam conceito 4, o conceito modal.
Dos demais, um curso recebeu o conceito 2, dois obtiveram o conceito 3, e a um curso foi atribuído o conceito 5.
Apenas um curso ficou sem conceito ( SC ) nessa categoria e região.
Na região Centro-Oeste, 12 dos 23 cursos eram de Universidades.
Nesse tipo de organização, nenhum curso ficou sem conceito ( SC ) ou com conceito 1.
O conceito modal foi 3, atribuído a sete cursos, seguido dos conceitos 2 e 4 atribuídos a dois cursos cada.
Apenas um curso ficou com conceito 5.
Os Centros Universitários da região Centro-Oeste contaram com quatro cursos.
Esses obtiveram conceitos 1 e 3 ( um curso cada ) e 4, conceito modal ( dois cursos ).
Nenhum curso ficou sem conceito ( SC ) ou recebeu os conceitos 2 ou 5.
Dos sete cursos oferecidos em Faculdades na região Centro-Oeste, dois ficaram alocados no conceito 2, e quatro cursos obtiveram conceito 3, conceito modal.
Apenas um curso ficou sem conceito ( SC ) enquanto nenhum outro curso foi avaliado com os demais conceitos.
O Gráfico 5.3 apresenta a distribuição do Conceito Enade segundo a Organização Acadêmica da IES.
Os cursos em Universidades ( linha azul ) apresentam uma poligonal mais à direita do que as demais, denotando uma melhor distribuição de Conceitos Enade.
No outro extremo, tem -se os cursos oferecidos em Centros Universitários ( linha verde ) com a poligonal mais à esquerda e, consequentemente, uma distribuição de Conceitos Enade com menores valores.
IES de Faculdades ( linha vermelha ) apresentam uma situação intermediária.

% matched lemmas: observação
\end{textsample}
